% Compile with XeLaTeX; TeX Live with ctex and Fandol fonts required.
\documentclass[UTF8,a4paper,11pt,fontset=fandol]{ctexart}
\usepackage[margin=23mm,headheight=14pt]{geometry}
\usepackage{amsmath,amssymb}
\usepackage{microtype}
\usepackage{xcolor}
\usepackage{fancyhdr}
\usepackage[unicode,hidelinks]{hyperref}
\definecolor{accent}{HTML}{315B75}
\ctexset{section={format=\Large\bfseries\color{accent}},subsection={format=\large\bfseries}}
\setlength{\parindent}{0pt}
\setlength{\parskip}{5pt}
\setlength{\emergencystretch}{2em}
\allowdisplaybreaks[2]
\pagestyle{fancy}
\fancyhf{}
\fancyhead[L]{\small 高等数学 A（上）期中试题与详解}
\fancyhead[R]{\small 2025--2026 学年}
\fancyfoot[C]{\thepage}
\title{\textbf{从极限到中值定理}\\[5pt]\Large 2025--2026 高等数学 A（上）期中试题与详解}
\author{整理：zeta2}
\date{2026 年 11 月 10 日}
\begin{document}
\maketitle
本文依据北京邮电大学 2025—2026 学年第一学期《高等数学 A（上）》期中试题及配套答案整理，保留原题编号与分值，并补充计算过程和证明细节。

\textbf{整理：zeta2}

试卷满分 100 分，包括填空题、选择题、极限计算、参数方程求导及中值定理证明。



\section*{一、填空题（共 30 分，每小题 5 分）}

\subsection*{1. 反三角函数的定义域}

\textbf{题目：} 已知 $f(x)=\arccos(1-x^2)$，求定义域。

\textbf{解：} 反余弦函数要求自变量属于 $[-1,1]$，因此
\[
-1\le 1-x^2\le 1
\quad\Longleftrightarrow\quad 0\le x^2\le 2.
\]
故定义域为
\[
\boxed{[-\sqrt{2},\sqrt{2}]}.
\]

\subsection*{2. 重要极限}

\textbf{题目：} 求
\[
\lim_{x\to\infty}\left(\frac{1+x}{x}\right)^{2x}.
\]

\textbf{解：} 利用重要极限，得
\[
\lim_{x\to\infty}\left(1+\frac1x\right)^{2x}
=\left[\lim_{x\to\infty}\left(1+\frac1x\right)^x\right]^2
=\boxed{e^2}.
\]

\subsection*{3. 无穷小的阶}

\textbf{题目：} 当 $x\to0$ 时，$\sin x-\tan x$ 与 $x^n$ 是同阶无穷小，求 $n$。

\textbf{解：} 因为
\[
\sin x-\tan x
=\frac{\sin x(\cos x-1)}{\cos x},
\]
所以
\[
\lim_{x\to0}\frac{\sin x-\tan x}{x^3}
=\lim_{x\to0}\frac{\sin x}{x}
\frac{\cos x-1}{x^2}\frac1{\cos x}
=-\frac12.
\]
极限为有限非零常数，故 $\boxed{n=3}$。

\subsection*{4. 复合函数求导}

\textbf{题目：} 已知 $y=\arctan(x^2+1)$，求 $\left.\dfrac{\mathrm{d}y}{\mathrm{d}x}\right|_{x=0}$。

\textbf{解：}
\[
y'=\frac{2x}{1+(x^2+1)^2},
\qquad \boxed{y'(0)=0}.
\]

\subsection*{5. 对数求导法}

\textbf{题目：} 设 $y=x^{\sin x}$，$x>0$，求 $y'$。

\textbf{解：} 两边取对数，有 $\ln y=\sin x\ln x$。求导得
\[
\frac{y'}y=\cos x\ln x+\frac{\sin x}{x}.
\]
因此
\[
\boxed{y'=x^{\sin x}\left(\frac{\sin x}{x}+\cos x\ln x\right)}.
\]

\subsection*{6. 分段函数的可导性}

\textbf{题目：} 函数
\[
f(x)=\begin{cases}
ax^2+b,&x\ge1,\\
x\cos\left(\dfrac{\pi x}{2}\right),&x<1
\end{cases}
\]
在 $x=1$ 处可导，求 $a,b$。

\textbf{解：} 可导必连续，故
\[
a+b=\lim_{x\to1^-}x\cos\left(\frac{\pi x}{2}\right)=0.
\]
在满足连续性的条件下，左右导数分别为
\[
f'_-(1)=\left[\cos\left(\frac{\pi x}{2}\right)
-\frac{\pi x}{2}\sin\left(\frac{\pi x}{2}\right)\right]_{x=1}
=-\frac\pi2,
\qquad f'_+(1)=2a.
\]
令二者相等，得
\[
\boxed{a=-\frac\pi4,\qquad b=\frac\pi4}.
\]

\section*{二、选择题（共 30 分，每小题 5 分）}

\subsection*{1. 幂指函数的极限}

\textbf{题目：} 求 $\displaystyle\lim_{x\to0}(\cos^2x)^{1/x^2}$。

A. $e^{-1/2}$；B. $1$；C. $e$；D. $e^{-1}$。

\textbf{解：} 底数在零点附近为正。先计算对数的极限：
\[
\lim_{x\to0}\frac{\ln(\cos^2x)}{x^2}
=\lim_{x\to0}\frac{\ln(1-\sin^2x)}{-\sin^2x}
\frac{-\sin^2x}{x^2}=-1.
\]
所以原极限为 $e^{-1}$，选 \textbf{D}。

\subsection*{2. 导数定义与复合增量}

\textbf{题目：} 设 $f'(0)=4$，求
\[
\lim_{x\to0}\frac{f(3\sin x)-f(2\arctan x)}x.
\]

A. $\dfrac12$；B. $2$；C. $4$；D. $\dfrac14$。

\textbf{解：} 由 $f$ 在零点可导，有
\[
f(u)=f(0)+4u+o(u)\qquad(u\to0).
\]
分别令 $u=3\sin x$ 和 $u=2\arctan x$，得
\[
f(3\sin x)-f(2\arctan x)
=4(3\sin x-2\arctan x)+o(x).
\]
因此原极限为 $4(3-2)=4$，选 \textbf{C}。

\subsection*{3. 间断点的分类}

\textbf{题目：} 设
\[
f(x)=\frac{\sin x}{|\tan x|},
\]
判断 $x=0$ 是哪一类间断点。

A. 可去间断点；B. 跳跃间断点；C. 第二类间断点；D. 无法判断。

\textbf{解：} 注意分母中的绝对值。在零点的充分小去心邻域内，$\cos x>0$，故
\[
f(x)=\frac{\sin x}{|\sin x|}\cos x.
\]
于是
\[
\lim_{x\to0^-}f(x)=-1,\qquad
\lim_{x\to0^+}f(x)=1.
\]
左右极限均存在但不相等，故为跳跃间断点，选 \textbf{B}。

\subsection*{4. 函数增量的线性主部}

\textbf{题目：} 设 $f$ 可导。当自变量在 $x=1$ 处取得增量 $\Delta x=0.1$ 时，相应函数增量 $\Delta y$ 的线性主部为 $0.5$，求 $f'(1)$。

A. $1$；B. $0.1$；C. $5$；D. $0.5$。

\textbf{解：} 线性主部为 $f'(1)\Delta x$，所以
\[
f'(1)\cdot0.1=0.5\quad\Longrightarrow\quad f'(1)=5.
\]
选 \textbf{C}。

\subsection*{5. 隐函数求导}

\textbf{题目：} 函数 $y=y(x)$ 由方程
\[
\ln(x^2+y)-x^3y-\sin x=0
\]
确定，求 $\left.\dfrac{\mathrm{d}y}{\mathrm{d}x}\right|_{x=0}$。

A. $0$；B. $1$；C. $-1$；D. $\dfrac12$。

\textbf{解：} 先代入 $x=0$，得到 $\ln y=0$，即 $y(0)=1$。对原方程两边求导：
\[
\frac{2x+y'}{x^2+y}-3x^2y-x^3y'-\cos x=0.
\]
再代入 $(x,y)=(0,1)$，得到 $y'(0)-1=0$，故选 \textbf{B}。

\subsection*{6. 高阶导数}

\textbf{题目：} 设 $f(x)=\ln(1-x)$，求 $f^{(2025)}(0)$。

A. $-2024!$；B. $2023!$；C. $2024!$；D. $2025!$。

\textbf{解：} 逐次求导可归纳得到
\[
f^{(n)}(x)=-\frac{(n-1)!}{(1-x)^n},\qquad n\ge1.
\]
所以 $f^{(2025)}(0)=-2024!$，选 \textbf{A}。

\textbf{选择题答案汇总：} D、C、B、C、B、A。

\section*{三、极限计算（15 分）}

\subsection*{1. 三角函数之差的极限（7 分）}

\textbf{题目：} 求
\[
\lim_{x\to0}\frac{\tan x-x}{4x^3}.
\]

\textbf{解：} 分子、分母均趋于零，使用洛必达法则：
\[
\begin{aligned}
\lim_{x\to0}\frac{\tan x-x}{4x^3}
&=\lim_{x\to0}\frac{\sec^2x-1}{12x^2}\\
&=\lim_{x\to0}\frac{\sin^2x}{12x^2\cos^2x}\\
&=\boxed{\frac1{12}}.
\end{aligned}
\]

\subsection*{2. 幂平均的极限（8 分）}

\textbf{题目：} 设 $a_1,a_2,\ldots,a_n>0$，求
\[
\lim_{x\to0}\left(\frac{a_1^x+a_2^x+\cdots+a_n^x}{n}\right)^{1/x}.
\]

\textbf{解：} 记
\[
A(x)=\frac1n\sum_{i=1}^n a_i^x.
\]
则 $A(x)>0$，且 $A(0)=1$。先求对数的极限，由洛必达法则，
\[
\begin{aligned}
\lim_{x\to0}\frac{\ln A(x)}x
&=\lim_{x\to0}\frac{A'(x)}{A(x)}\\
&=\lim_{x\to0}\frac{\sum_{i=1}^n a_i^x\ln a_i}{\sum_{i=1}^n a_i^x}\\
&=\frac1n\sum_{i=1}^n\ln a_i.
\end{aligned}
\]
由指数函数的连续性，原极限为
\[
\boxed{\exp\left(\frac1n\sum_{i=1}^n\ln a_i\right)
=\sqrt[n]{a_1a_2\cdots a_n}}.
\]
这一写法也涵盖所有 $a_i=1$ 的情形，无需额外除以 $A(x)-1$。

\section*{四、参数方程求导（15 分）}

\textbf{题目：} 设参数方程
\[
\begin{cases}
x=e^t+t,\\
e^t+e^y=2
\end{cases}
\]
确定函数 $y=y(x)$，求 $x=1$ 处的切线、法线方程，以及 $\left.\dfrac{\mathrm{d}^2y}{\mathrm{d}x^2}\right|_{x=1}$。

\subsection*{1. 切线与法线（9 分）}

由 $e^t+e^y=2$ 可得
\[
y=\ln(2-e^t),\qquad t<\ln2.
\]
又因为 $\dfrac{\mathrm{d}x}{\mathrm{d}t}=e^t+1>0$，$x$ 随 $t$ 严格递增。当 $x=1$ 时，唯一对应 $t=0$，进而 $y=0$。

参数方程的一阶导数为
\[
\frac{\mathrm{d}y}{\mathrm{d}x}
=\frac{\mathrm{d}y/\mathrm{d}t}{\mathrm{d}x/\mathrm{d}t}
=-\frac{e^t}{(2-e^t)(e^t+1)}.
\]
因此在 $(1,0)$ 处，切线斜率为 $-\dfrac12$，法线斜率为 $2$。两条直线分别为
\[
\boxed{y=-\frac12(x-1)}\quad\text{（切线）},
\]
\[
\boxed{y=2(x-1)}\quad\text{（法线）}.
\]

\subsection*{2. 二阶导数（6 分）}

记
\[
p(t)=\frac{\mathrm{d}y}{\mathrm{d}x}
=-\frac{e^t}{(2-e^t)(e^t+1)}.
\]
对 $t$ 求导，可整理为
\[
p'(t)=p(t)\left(1+\frac{e^t}{2-e^t}-\frac{e^t}{e^t+1}\right).
\]
于是
\[
p(0)=-\frac12,\qquad
p'(0)=-\frac12\left(1+1-\frac12\right)=-\frac34.
\]
注意二阶导数还需除以 $x'(t)$：
\[
\frac{\mathrm{d}^2y}{\mathrm{d}x^2}
=\frac{p'(t)}{x'(t)}.
\]
故
\[
\boxed{\left.\frac{\mathrm{d}^2y}{\mathrm{d}x^2}\right|_{x=1}
=\frac{-3/4}{2}=-\frac38}.
\]

\section*{五、中值定理证明（10 分）}

\textbf{题目：} 设 $f$ 在 $[0,\pi]$ 上连续，在 $(0,\pi)$ 内可导，且对任意 $x\in(0,\pi)$ 都有 $f'(x)\ne0$。证明下列结论。

\subsection*{1. 拉格朗日中值定理（4 分）}

\textbf{待证：} 存在 $c\in(0,\pi)$，使得
\[
\frac{f(\pi)-f(0)}\pi=f'(c).
\]

\textbf{证明：} 由题设，$f$ 满足闭区间 $[0,\pi]$ 上的拉格朗日中值定理条件。因此存在 $c\in(0,\pi)$，使得
\[
f(\pi)-f(0)=f'(c)(\pi-0).
\]
两边除以 $\pi$ 即得结论。

\subsection*{2. 构造辅助函数与罗尔定理（6 分）}

\textbf{待证：} 存在 $\xi\in(0,\pi)$，使得
\[
\frac{f(\xi)}{f'(\xi)}=-\tan\xi.
\]

\textbf{证明：} 令
\[
g(x)=f(x)\sin x.
\]
则 $g$ 在 $[0,\pi]$ 上连续，在 $(0,\pi)$ 内可导，并且
\[
g(0)=f(0)\sin0=0,\qquad g(\pi)=f(\pi)\sin\pi=0.
\]
由罗尔定理，存在 $\xi\in(0,\pi)$，使得
\[
g'(\xi)=f'(\xi)\sin\xi+f(\xi)\cos\xi=0.
\]
还需说明 $\cos\xi\ne0$。若 $\cos\xi=0$，则因 $\xi\in(0,\pi)$，必有 $\xi=\dfrac\pi2$。此时上式变为 $f'(\xi)=0$，与题设矛盾。

故 $f'(\xi)\cos\xi\ne0$，可以将等式两边除以它，得到
\[
\boxed{\frac{f(\xi)}{f'(\xi)}=-\frac{\sin\xi}{\cos\xi}=-\tan\xi}.
\]
证毕。

\section*{整理说明与资料}

题目依据《25-26高数上-期中终稿.pdf》，答案依据《25-26高数上-期中答案.pdf》。本文为学习整理版；题目及参考答案来源于所附试卷资料，解题过程作了补充和排版整理。

整理时保留了选择题第 3 题分母的绝对值，修正了参考答案中个别排版笔误，并补充了最后一题中 $\cos\xi\ne0$ 的论证。

\end{document}
